% Motivation and learning outcomes.
% The outcomes are the outcomes of Day 4 in the syllabus.
\begin{frame}{Why this day matters}
  \begin{itemize}
    \item The model of Day 3 uses 32 bits for each weight and each activation.
          A microcontroller pays for each bit with flash, RAM, time, and
          energy.
    \item Quantization stores the same model with 8 bits for each value. The
          model becomes four times smaller, and integer kernels run faster.
    \item The cost is a rounding error. You must know where the error comes
          from and how to keep it small.
    \item In the lab you quantize tensors by hand, quantize a small CNN, and
          compare the float model and the int8 model on the XIAOML Kit.
  \end{itemize}
\end{frame}

\begin{frame}{Learning outcomes}
  After this day you can:
  \begin{itemize}
    \item Derive the scale and the zero point of an affine quantizer.
    \item Apply post-training quantization with a representative dataset.
    \item Apply quantization-aware training.
    \item Explain per-tensor and per-channel quantization.
    \item Measure the effect of int8 quantization on size, accuracy, and
          latency.
  \end{itemize}
\end{frame}
